\section{Dual-Scope Sequential Pattern Mining}\label{sec:sequences}

\subsection{From events to sequences}
Each patient's clinical history is turned into an ordered sequence of typed
items. An item is a clinical event prefixed by its type --- \texttt{C:} for a
condition, \texttt{M:} for a medication, \texttt{P:} for a procedure and
\texttt{CP:} for a care plan --- so that, e.g., \texttt{C:diabetes} and a
similarly named procedure remain distinct. Events are sorted chronologically;
events on the same day are ordered by a clinically motivated type rank
($C\!<\!M\!<\!P\!<\!CP$). Consecutive duplicates (such as repeated refills) are
collapsed, a curated blacklist removes administrative noise, only the most recent
window of history is kept, and the per-patient sequence length is capped. A
typical sequence is
\[
  \texttt{C:diabetes} \rightarrow \texttt{M:metformin} \rightarrow
  \texttt{P:foot\_exam} \rightarrow \texttt{C:hypertension}.
\]

\subsection{Dual-scope mining}
We apply PrefixSpan~\cite{Pei2004PrefixSpan} at two scopes. \emph{Global} mining
runs over the whole cohort and yields population-wide trajectories.
\emph{Per-cluster} mining runs PrefixSpan independently over the patients of each
cluster produced in Section~\ref{sec:clustering} (the outlier group included),
yielding trajectories characteristic of that subgroup. Comparing the two scopes
is the key analytical move: a trajectory may be dominant inside one cluster yet
rare in the population, which signals a group-specific care progression rather
than a generic one. Mined patterns are stored indexed by cluster label, so the
front end can request either the global set or any cluster's set, and each
pattern is linked to the patients that exhibit it.

\subsection{Ranking and filtering}
Frequency alone over-ranks patterns built from individually common items. We
score each pattern $P=\langle i_1,\dots,i_n\rangle$ by its \emph{geometric-mean
lift}, comparing the pattern's observed support to what independence would
predict:
\begin{equation}
  \mathrm{GML}(P) \;=\;
  \frac{\mathrm{supp}(P)/N}
       {\left(\prod_{k=1}^{n}\mathrm{supp}(i_k)/N\right)^{1/n}},
  \label{eq:gml}
\end{equation}
computed in log-space for numerical stability, where $N$ is the number of
patients in scope. A pattern whose items rarely co-occur scores higher than one
made of ubiquitous items. We further multiply the score by a
\emph{type-diversity bonus} that rewards patterns spanning several event types
(condition $\rightarrow$ medication $\rightarrow$ procedure), since these tell a
richer clinical story than repetitions of a single type. Mining requests closed
patterns of length 3--6 above a $1\%$ support floor; a larger candidate pool is
over-mined and then reduced by the ranking above plus a set of filters: removal
of runs of three same-type items, an explicit closedness check, Jaccard-based
near-duplicate suppression, and per-prefix and per-suffix caps that prevent a
single dominant condition from monopolizing the result list. The top-ranked
patterns per scope are persisted for presentation.
